% Einheit 1 von 3 – Kreisumfang (bank/kreis/e1.jsonl, zusammenbau v0.1)
%% TODO Zweigzeile Teil 1: was hier gelernt wird (Satz in Schülersprache aus dem Einheitstitel) – Einheit 1
\einheitenkopf[e1]{Einheit 1 von 3 · Kreisumfang}
\zweigzeile{ab Kl. 7, je nach Buch bis Kl. 8 · P10}
\setcounter{aufgabe}{8}

%% TODO Titel: Ich-kann-Satz fehlt in der Bank (Platzhalter: Kettenname) – Nr. 9
%% TODO Anweisung: ein Satz über der ersten Teilaufgabe fehlt in der Bank – Nr. 9
\begin{aufgabe}{Rand oder Fläche?}
\begin{teile}
\teil Ein runder Tisch bekommt rundherum eine Kante aus Holzleiste. Rand oder Fläche? Kreuze an.\\ \kreuz{Umfang}\\ \kreuz{Fläche}
\end{teile}
\end{aufgabe}

%% TODO Titel: Ich-kann-Satz fehlt in der Bank (Platzhalter: Kettenname) – Nr. 10
%% TODO Anweisung: ein Satz über der ersten Teilaufgabe fehlt in der Bank – Nr. 10
\begin{aufgabe}{Umfang}
\swz{Ist die beschriftete Strecke ein Radius oder ein Durchmesser? Gib auch die andere Länge an.}{\begin{kreis}[2] \mittelpunkt{M} \durchmesser{20}{$9$ cm} \end{kreis}}
%% TODO Darstellung für schwach fehlt in der Bank (kreis-e1-k2-s1-v1; Abschnitt „Für schwache Schüler“)
\swz{Kreis mit $d = 6$ cm – wie lang ist der Umfang? Rechne mit der $\pi$-Taste und runde auf eine Stelle.}{}
%% TODO Darstellung für schwach fehlt in der Bank (kreis-e1-k2-s1-v2; Abschnitt „Für schwache Schüler“)
\swz{Kreis mit $d = 11$ cm – wie lang ist der Umfang? Rechne mit der $\pi$-Taste und runde auf eine Stelle.}{}
%% TODO Darstellung für schwach fehlt in der Bank (kreis-e1-k2-s1-v3; Abschnitt „Für schwache Schüler“)
\swz{Kreis mit $d = 3$ m – wie lang ist der Umfang? Rechne mit der $\pi$-Taste und runde auf eine Stelle.}{}
%% TODO Darstellung für schwach fehlt in der Bank (kreis-e1-k2-s1-v4; Abschnitt „Für schwache Schüler“)
\swz{Kreis mit $d = 13$ cm – wie lang ist der Umfang? Rechne mit der $\pi$-Taste und runde auf eine Stelle.}{}
%% TODO Darstellung für schwach fehlt in der Bank (kreis-e1-k2-s1-v5; Abschnitt „Für schwache Schüler“)
\swz{Kreis mit $d = 17$ mm – wie lang ist der Umfang? Rechne mit der $\pi$-Taste und runde auf eine Stelle.}{}
\swfrage{Was bleibt gleich, was ändert sich?}
%% TODO Darstellung für schwach fehlt in der Bank (kreis-e1-k2-s2-v1; Abschnitt „Für schwache Schüler“)
\swz{Kreis mit $r = 4$ cm – wie lang ist der Umfang? Rechne mit der $\pi$-Taste und runde auf eine Stelle.}{}
%% TODO Darstellung für schwach fehlt in der Bank (kreis-e1-k2-s3-v1; Abschnitt „Für schwache Schüler“)
\swz{Kreis mit $d = 4{,}6$ cm – wie lang ist der Umfang? Rechne mit der $\pi$-Taste und runde auf zwei Stellen.}{}
%% TODO Darstellung für schwach fehlt in der Bank (kreis-e1-k2-s4-v1; Abschnitt „Für schwache Schüler“)
\swz{Ein Kreis hat den Umfang $u = 47$ cm. Wie groß ist sein Durchmesser? Rechne mit der $\pi$-Taste und runde auf eine Stelle.}{}
%% TODO Darstellung für schwach fehlt in der Bank (kreis-e1-k2-s5-v1; Abschnitt „Für schwache Schüler“)
\swz{Ein Kreis hat den Umfang $u = 75$ cm. Wie groß ist sein Radius? Rechne mit der $\pi$-Taste und runde auf eine Stelle.}{}
%% TODO Darstellung für schwach fehlt in der Bank (kreis-e1-k2-s6-v1; Abschnitt „Für schwache Schüler“)
\swz{Ein Weg führt im Halbkreis um einen Teich, der Durchmesser des Halbkreises ist $d = 12$ m. Wie lang ist der Weg? Rechne mit der $\pi$-Taste und runde auf eine Stelle.}{}
%% TODO Darstellung für schwach fehlt in der Bank (kreis-e1-k2-s7-v1; Abschnitt „Für schwache Schüler“)
\swz{Ein Fahrradrad hat den Durchmesser $66$ cm. Wie weit fährt das Rad bei $250$ Umdrehungen? Gib das Ergebnis in Meter an, auf eine Stelle gerundet.}{}
\swz[3]{Ein Stifteköcher ohne Deckel hat die Form eines Zylinders mit dem Radius $3{,}6$ cm und der Höhe $11$ cm. Ein rechteckiges Papier soll genau einmal um die Seitenfläche reichen. Berechne die Länge und die Breite des Rechtecks; runde auf eine Stelle.}{}
\end{aufgabe}

%% TODO Titel: Ich-kann-Satz fehlt in der Bank (Platzhalter: Kettenname) – Nr. 11
%% TODO Anweisung: ein Satz über der ersten Teilaufgabe fehlt in der Bank – Nr. 11
\begin{aufgabe}{Radius, Durchmesser, Mittelpunkt in einer Figur benennen und einzeichnen}
\swa{Zeichne in den Kreis einen Durchmesser ein, der durch den Punkt $P$ geht, und beschrifte ihn mit $d$.}{\begin{kreis}[2] \mittelpunkt{M} \kreispunkt{30}{P} \end{kreis}}
\end{aufgabe}

%% TODO Titel: Ich-kann-Satz fehlt in der Bank (Platzhalter: Kettenname) – Nr. 12
%% TODO Anweisung: ein Satz über der ersten Teilaufgabe fehlt in der Bank – Nr. 12
\begin{aufgabe}{Kreis mit gegebenem Radius oder Durchmesser zeichnen}
\swa{Zeichne einen Kreis mit dem Radius $3$ cm und markiere seinen Mittelpunkt mit $M$.}{\rechenplatz[halb]{6}}
\end{aufgabe}

%% TODO Titel: Ich-kann-Satz fehlt in der Bank (Platzhalter: Kettenname) – Nr. 13
%% TODO Anweisung: ein Satz über der ersten Teilaufgabe fehlt in der Bank – Nr. 13
\begin{aufgabe}{Umfang messen und u : d bilden (Entdeckung von π)}
%% TODO Darstellung für schwach fehlt in der Bank (kreis-e1-k5-s1-v1; Abschnitt „Für schwache Schüler“)
\swz{An einem Becher wurde gemessen: Umfang $u = 22$ cm, Durchmesser $d = 7$ cm. Berechne $u : d$ und runde auf zwei Stellen.}{}
\end{aufgabe}

%% TODO Titel: Ich-kann-Satz fehlt in der Bank (Platzhalter: Kettenname) – Nr. 14
%% TODO Anweisung: ein Satz über der ersten Teilaufgabe fehlt in der Bank – Nr. 14
\begin{aufgabe}{Umfang – Fehler finden, Begründen, Anwendung, Darstellungswechsel}
\swz[3]{Ein Kreis hat den Durchmesser $d = 18$ cm. Ben rechnet den Umfang: \rechnung{u &= 2 \cdot \pi \cdot 18 \\ u &\approx 113{,}1 \text{ cm}} Wo ist der Fehler?}{}
\swfrage{Begründe, warum man bei jedem Kreis für $u : d$ ungefähr dasselbe Ergebnis erhält.}
\swz[3]{Eine Försterin misst den Umfang einer Eiche mit dem Maßband: $2{,}83$ m. Ab $75$ cm Durchmesser steht ein Baum hier unter Schutz. Steht die Eiche unter Schutz?}{}
\swa{Ergänze die Tabelle. Runde auf eine Stelle.}{\wertetabelle{$d$ in cm}{$u$ in cm}{1,2,3,6}}
\end{aufgabe}

\uebersichtskasten{Umfang: Durchmesser mal π (π ≈ 3,14) – der Umfang ist gut dreimal so lang wie der Durchmesser. \\ d = 5 cm: u = π · 5 ≈ 15,7 cm \quad r = 2 cm: d = 4 cm, u = π · 4 ≈ 12,6 cm \\ Rückwärts: d = u : π. \quad u = 31,4 cm: d = 31,4 : π ≈ 10 cm}
